\section{Domain-Selective Verbalized Rule Injection}
\label{sec:verbalized}

Passive the refutation agent (the refutation agent without any rulebook) re-reads the normal
reference images on every test query but never \emph{accumulates}
anything from them: the same ref set produces the same ad-hoc
reasoning trace a thousand times. This section asks whether we can
externalise domain knowledge --- either from the refs themselves or
from a small oracle-labelled dev budget --- into a persistent
natural-language rulebook that the agent consults at test time.
Our empirical answer is \emph{domain-selective yes}: a
one-sentence task anchor helps everywhere, oracle-grounded cluster
rules help on a majority of domains but not uniformly, and
self-supervised invariant extraction helps on structural/logical
domains while hurting on medical-pathology ones. The macro
point-estimates move by $+2$--$3$ pp but the paired 95\% CIs
include zero, so we foreground heterogeneity rather than claim a
universal ``self-evolution'' improvement.

\subsection{Four learning regimes}
\label{ssec:verb_regimes}

We frame the design as a ladder of four regimes, ordered by how much
supervision the system touches and how aggressive the textualisation
is:

\begin{description}
\item[\textbf{Passive the refutation agent} (training-free)] — no rulebook. Baseline.
\item[\textbf{+Anchor} (zero oracle)] — a one-sentence task anchor
  per domain (concatenating the domain's \texttt{anomaly\_type} and
  free-text description from the benchmark spec) is prepended to the
  agent's turn-$1$ system prompt. No rulebook content beyond the
  anchor. This is the \S\ref{sec:experiments} Table~\ref{tab:descriptor_ablation}
  descriptor finding, applied at the agent (not just Direct) level.
\item[\textbf{+L1 invariants} (self-supervised, zero oracle)] — a
  reflector VLM reads $n{=}8$ normal reference images per
  (domain, class) and emits a small list of \emph{invariants} ---
  concrete visual predicates that hold in every ref. Invariant types
  are restricted to six categories: \texttt{count}, \texttt{symmetry},
  \texttt{spatial\_layout}, \texttt{color\_palette}, \texttt{texture},
  \texttt{structural}. Anomaly at test time is defined as
  \emph{invariant violation}. Crucially, the reflector is forbidden
  from listing ``hypothesised anomaly modes'' (a design that an
  earlier iteration attempted but that collapsed into generic
  VLM-prior recitation; Appendix~\ref{app:verbalized_v1_pilot}).
  If no invariant can be verified across all refs the list is empty.
\item[\textbf{+L2 clusters} (semi-supervised, $K{=}10$ oracle/domain)]
  — run Passive the refutation agent on the dev split, form a balanced batch of
  $K/2$ false negatives (gt=anomaly, score{$<$}0.5) and $K/2$ false
  positives (gt=normal, score{$\ge$}0.5), reveal their manifest
  labels, and ask a reflector VLM --- seeing the \emph{whole batch}
  at once --- to propose $1$--$3$ corrective rules per side that
  capture common patterns across multiple items. Single-item patterns
  are explicitly forbidden. \texttt{normal\_tolerance} rules must
  name a specific visual feature, not blanket-suppress an anomaly
  class. This per-batch clustering was a crucial difference from a
  per-item predecessor that over-generated false \texttt{normal\_tolerance}
  rules; see \S\ref{ssec:verb_method} for protocol details.
\item[\textbf{+L1+L2} (stacked, $K{=}10$ oracle/domain)] — offline
  concatenation of the L1 invariant list and the L2 cluster rules,
  deduplicated; retrieval is carried out against the union store.
\end{description}

All four rule-consuming regimes share a retrieval-augmented
injection mechanism: the agent is handed only the top-$k{=}3$ rules
that match the query's (domain, class) metadata, filtered and ranked
by (i) priority (L2 false-negative corrective first, then L1
invariants, then L2 false-positive corrective), (ii) natural order
within tier, with the constraint that at least one L1 invariant is
retained when available. This keeps the per-query payload under
$\sim 100$ tokens, in contrast to the full rulebook dump that an
earlier iteration injected ($\sim 250$ tokens per query; it incurred
$10$--$24\%$ JSON parse-failure rates, against $\le 3\%$ under the
RAG-based injection used here).

\subsection{Method}
\label{ssec:verb_method}

\paragraph{Dev pre-split.} Each domain's 40-item dev split is
partitioned deterministically into a 20-item \emph{selection} subset
$\mathcal{D}_{\mathrm{sel}}$ (oracle pool) and a 20-item
\emph{validation} subset $\mathcal{D}_{\mathrm{val}}$ (untouched by
any reflector). The split is fixed by \texttt{seed}$=0$ throughout
this paper.

\paragraph{Granularity.} Multi-class domains (D1 industrial,
D2 retail, D3 complex industrial, D4 infrastructure,
D5 logical, D6 3D industrial) receive \emph{per-class}
rulebooks; single-class domains (D7 remote sensing, D8 dermatology,
D9 brain MRI, D10 liver CT, D11 GI endoscopy, D12 road safety) receive
per-domain rulebooks. Class identity at test time is read from the
manifest (no oracle prediction).

\paragraph{L1 invariant extraction.} For each (domain, class) the
reflector sees $8$ NORMAL references and the domain task anchor, and
returns a JSON object with up to $6$ typed invariants plus a
\texttt{confidence} self-assessment (high / medium / low). Empty
invariant lists are allowed and occur on weak-prior domains (D10
liver CT, D6 Real3D-AD for some classes): the reflector is
instructed to return an empty list rather than fabricate generic
rules.

\paragraph{L2 batch reflection.} The oracle budget is
\textbf{$K{=}10$ revealed labels per domain, shared across classes}
(not per class). For each domain we run Passive the refutation agent on
$\mathcal{D}_{\mathrm{sel}}$, form the balanced $K/2$+$K/2$ batch
across all classes of that domain, and call the reflector once with
the full batch (all $K$ query images, all agent trajectories, all
revealed labels). On single-class domains (D7--D12) this produces
one batch covering the whole domain; on multi-class domains (D1--D6)
classes receive uneven coverage depending on which classes contain
the hardest FN/FP cases, and we report per-class coverage counts in
Appendix~\ref{app:verbalized_coverage}. The reflector proposes
$1$--$3$ rules per side (FN cluster and FP cluster), each covering
at least two items, and classifies each as
\texttt{anomaly\_mode\_extension} / \texttt{normal\_tolerance} /
\texttt{rule} with a strict instruction that \texttt{normal\_tolerance}
must name a specific pattern rather than blanket-suppress. This
yielded $2$--$6$ rules per (domain, class bucket), rather than the
per-item $10$ rules a predecessor produced.

\paragraph{Test-time injection.} The agent's unchanged the refutation agent system
prompt is concatenated with the anchor. The top-$3$ retrieved rules
(see \S\ref{ssec:verb_regimes}) are appended to the turn-$1$ user
message as a brief markdown list tagged by rule type. No test items
enter the rulebook-construction pipeline at any stage.

\subsection{Results}
\label{ssec:verb_exp}

All results below use Qwen3.5-VL-27B as the backbone, max 3 turns
per query, temperature 0, and the strict dev-only protocol above.
The baseline is \emph{Passive the refutation agent on manifests\_v2 test}
($n{=}1{,}418$, 12 domains) — a fresh run that does not reuse §4's
/.5 agent numbers because manifests\_v2 resamples items on most
domains and §4's main-table agent architecture differs.

\begin{table}[t]
\centering
\caption{\textbf{Verbalized Self-Evolution — macro AUROC} on
CrossDomainVAD-11 manifests\_v2 test (12 domains, $n{=}1{,}418$,
Qwen3.5-VL-27B). Four rulebook regimes against the Passive the refutation agent
baseline. Bootstrap 95\% CIs on macro $\Delta$ via $1{,}000$
stratified paired resamples over domains.}
\label{tab:verb_main}
\small
\setlength{\tabcolsep}{5pt}
\begin{tabular}{lccccc}
\toprule
Variant & Oracle / domain & Payload (tok) & Macro AUROC & $\Delta$ vs Passive [95\% CI] & $\hat{p}(\Delta{>}0)$ \\
\midrule
Passive the refutation agent & $0$ & --- & $0.6685$ & --- & --- \\
\textbf{+ Anchor}   & $0$  & $\sim 20$ & $0.6885$ & $+2.00$ [$-0.67$, $+4.86$] & $0.934$ \\
\textbf{+ L1 inv}   & $0$  & $\sim 80$ & $0.6603$ & $-0.82$ [$-3.74$, $+2.17$] & $0.320$ \\
\textbf{+ L2 cls}   & $10$ & $\sim 80$ & $\mathbf{0.6958}$ & $\mathbf{+2.73}$ [$-0.46$, $+6.01$] & $\mathbf{0.957}$ \\
\textbf{+ L1+L2}    & $10$ & $\sim 90$ & $0.6708$ & $+0.23$ [$-2.70$, $+3.53$] & $0.566$ \\
\bottomrule
\end{tabular}
\end{table}

\paragraph{Finding 1: A one-sentence task anchor gives $+2.00$ pp
at zero oracle cost, consistent with \S\ref{sec:experiments}'s
descriptor finding at the agent level.} The anchor is taken
verbatim from \texttt{domain\_config.json}: it concatenates the
domain's \texttt{anomaly\_type} (e.g.\ ``building\_change'',
``melanoma'', ``polyp/ulcer'') with its free-text description.
Per-domain contributions concentrate on D11 GI endoscopy
($+13.6$ pp, anchor explicitly names ``polyp/ulcer''), D7 satellite
change detection ($+12.0$ pp), and D6 3D industrial ($+8.4$ pp).

\paragraph{Finding 2: Active verbalised learning via cluster-based
dev oracle gives the largest point-estimate gain ($+2.73$ pp) among
the four regimes.} The paired 95\% CI is $[-0.46, +6.01]$ --- it
includes zero by a thin margin, so we do not claim formal
significance at the conventional threshold. The bootstrap posterior
mass above zero, $\hat{p}(\Delta{>}0){=}0.957$, is reported
alongside the CI rather than as a significance test; we read this
as ``a reliable direction under the paired design, with a gain
magnitude comparable to Anchor but reached via a different
mechanism''. Replacing a per-item oracle reflector with a per-batch
clustering reflector, combined with balanced FN/FP selection on
$\mathcal{D}_{\mathrm{sel}}$, yields $2$--$6$ high-precision rules
per (domain, class). The largest per-domain gains are on D7
satellite ($+21.4$ pp), D4 SDNET cracks ($+15.4$ pp), and D9 brain
MRI ($+10.5$ pp). A predecessor with per-item reflection and
unconstrained oracle selection over-generated
\texttt{normal\_tolerance} rules that suppressed true anomalies on
medical domains; the cluster + constraint design in this paper
corrects that (full comparison in Appendix~\ref{app:verbalized_v1_pilot}).

\paragraph{Finding 3: Self-supervised invariants transfer only when
they describe the object class, not the image modality.} The L1
regime (net $-0.82$ pp) wins on domains where the refs encode
structural or combinatorial invariants whose \emph{violations
directly correspond to the anomaly signal}: D4 SDNET cracks
($+11.6$ pp; the invariant ``continuous unbroken concrete surface''
is violated by any crack), D3 VisA ($+3.5$ pp; count/placement
invariants are violated by missing or misplaced parts), D5
MVTec-LOCO ($+1.8$ pp; ``exactly $N$ components in grid'' is
violated by logical defects). On medical imaging domains the
reflector \emph{does} produce high-confidence invariants, but the
invariants describe the \emph{modality} rather than the object
class: on D9 brain MRI the extracted invariants are ``grayscale
palette, granular heterogeneous texture, central oval tissue
boundary'' --- all of which hold equally well on tumour slices ---
and similarly D10 liver CT gets ``black background, speckled
foreground'' which tumours share with normal tissue. Because the
invariants do not discriminate normal from anomaly, injecting them
as a rulebook \emph{adds noise} to the agent's reasoning rather
than signal ($-10.7$ / $-5.9$ / $-3.9$ pp on D8/D9/D10). The
practical consequence is that L1 is useful as a domain-selective
add-on (structural / logical / visible-defect domains), not as a
universal default. A cleaner test for ``invariant is
anomaly-discriminative'' is a promising extension: e.g.\ compute
the expected rulebook recall on an anomaly seed set, and suppress
rulebooks whose invariants hold on the seeds.

\paragraph{Finding 4: Stacking L1+L2 does not outperform L2 alone
on this setup.} Contrary to the original hypothesis, concatenating
L1 invariants with L2 cluster rules (net $+0.23$ pp) gives up the
L2-only win; the L1 invariants from the weaker domains dilute the
rulebook. We speculate that a confidence-gated stack (inject L1
only when \texttt{confidence}=high) would recover the L1 wins on
structural domains without the medical-domain drag; this is listed
as future work.

\subsection{Per-domain breakdown}
\label{ssec:verb_perdomain}

\begin{table}[t]
\centering
\caption{\textbf{Per-domain AUROC and $\Delta$ vs Passive the refutation agent} for
the four regimes on manifests\_v2 test, Qwen3.5-VL-27B. Largest
per-row gain in \textbf{bold}.}
\label{tab:verb_per_domain}
\small
\setlength{\tabcolsep}{3pt}
\begin{tabular}{llccccc}
\toprule
Dom & Source / defect type & Passive & +Anchor & +L1 inv & +L2 cls & +L1+L2 \\
\midrule
D1  & MVTec-AD / surface defect  & 0.918 & 0.897 ($-2.0$) & 0.895 ($-2.3$) & 0.851 ($-6.7$) & 0.874 ($-4.4$) \\
D2  & GoodsAD / packaging        & 0.598 & 0.574 ($-2.4$) & 0.577 ($-2.2$) & 0.639 ($+4.1$) & \textbf{0.672 ($+7.3$)} \\
D3  & VisA / fine-grained defect & 0.783 & 0.811 ($+2.8$) & \textbf{0.819 ($+3.5$)} & 0.748 ($-3.5$) & 0.812 ($+2.9$) \\
D4  & SDNET / concrete crack     & 0.581 & 0.555 ($-2.7$) & 0.697 ($+11.6$) & \textbf{0.735 ($+15.4$)} & 0.634 ($+5.3$) \\
D5  & MVTec-LOCO / logical       & 0.654 & 0.594 ($-6.0$) & \textbf{0.672 ($+1.8$)} & 0.664 ($+0.9$) & 0.600 ($-5.4$) \\
D6  & Real3D-AD / geometric      & 0.509 & \textbf{0.593 ($+8.4$)} & 0.469 ($-3.9$) & 0.507 ($-0.1$) & 0.578 ($+7.0$) \\
D7  & LEVIR / building change    & 0.614 & 0.734 ($+12.0$) & 0.644 ($+3.0$) & \textbf{0.828 ($+21.4$)} & 0.668 ($+5.4$) \\
D8  & DermaMNIST / melanoma      & 0.647 & \textbf{0.664 ($+1.7$)} & 0.540 ($-10.7$) & 0.549 ($-9.8$) & 0.522 ($-12.5$) \\
D9  & BraTS / tumor              & 0.720 & 0.722 ($+0.2$) & 0.661 ($-5.9$) & \textbf{0.825 ($+10.5$)} & 0.682 ($-3.8$) \\
D10 & BMAD-Liver / pathology     & 0.504 & 0.493 ($-1.1$) & 0.464 ($-3.9$) & 0.501 ($-0.3$) & 0.461 ($-4.3$) \\
D11 & HyperKvasir / polyp/ulcer  & 0.570 & \textbf{0.706 ($+13.6$)} & 0.568 ($-0.2$) & 0.581 ($+1.1$) & 0.624 ($+5.4$) \\
D12 & BDD+RoadAnomaly / obstacle & 0.924 & 0.921 ($-0.4$) & 0.917 ($-0.7$) & 0.922 ($-0.2$) & 0.923 ($-0.2$) \\
\midrule
\textbf{Macro} & & \textbf{0.669} & \textbf{0.689} ($+2.00$) & \textbf{0.660} ($-0.82$) & \textbf{0.696} ($+2.73$) & \textbf{0.671} ($+0.23$) \\
\bottomrule
\end{tabular}
\end{table}

Table~\ref{tab:verb_per_domain} gives the full per-domain picture.
The headline is that no single regime dominates across all 12
domains --- four different regimes each win on at least one domain
(Anchor: D6/D8/D11; L1 inv: D3/D5; L2 cls: D4/D7/D9; L1+L2: D2).
The per-domain winner-count, combined with the CI structure, is
consistent with the claim that verbalised rulebook injection is
\emph{useful and domain-selective}, rather than a universal
improvement.

\subsection{Limitations}
\label{ssec:verb_limitations}
The $P(\Delta{>}0){=}0.957$ for +L2 cls sits narrowly below the
conventional $95\%$ threshold; we report point estimates with full
CIs rather than claiming statistical significance at the standard
level. Three domains (D1 MVTec-AD, D10 liver CT, D12 road safety)
are largely unaffected by any rulebook regime in our experiments:
D1 has a very high Passive baseline ($0.918$) leaving little
headroom, D10 is near-random for Passive ($0.504$) suggesting the
backbone itself does not discriminate here, and D12 is ceiling
($0.924$). Rulebook injection does not help when Passive is already
saturated at either end. The L1 invariant regime has a clear
domain-type boundary between ``structural / logical / geometric''
domains (where it wins) and ``medical pathology / statistical
distribution'' domains (where it loses); confidence-gated injection
is left for future work, as is a turn-0 query-conditioned
retrieval step that would replace the static (domain, class)
metadata filter.
